\chapter{運動方程式を解く・中級}

中級といっても、中級の中でも下といったところになりましょう。
僕たちの手の中にある道具といえば、多項式関数や指数関数・対数関数の微積分が自由に行えるぐらい。
せめて三角関数さえあれば、もっと多くの微分方程式が解けるはずなのに…。



\section{空気抵抗ありの自由落下}

自由落下は、どちら向きに重力が働いているかによりますが、運動方程式は
\[
  m\ddot{\gamma}=(0,0,mg)
\]
で表されます。明らかに一次元での運動になるので、初めから$\gamma$は一次元の値をとるとし、
\[
  m\ddot{\gamma}=mg
\]
を考えましょう。

空気抵抗があるとき、それは$\gamma$の速さに比例する力が加わりそうという予感がします。
車の窓を開け、手を伸ばしてみると、車が速ければ速いほど、強い力を手に感じることができます\footnote{そんなことしちゃダメだぞ。}。
従って、自由落下の運動方程式を次のように変えるのが良さそうです。
\[
  m\ddot{\gamma}=mg-k\dot{\gamma}
\]
$k>0$は比例定数です。
$k$は、定性的な結果があっていれば、観測事実から後から決める実数です。
いったん、速度がわかればよいという観点から
\[
  v(t):=\dot{\gamma}(t)
\]
とおくと、運動方程式は
\[
  \frac{dv}{dt}=g-\frac{k}{m}v
\]
となります。

突然ですが、このように何かの未知関数の微分が絡んだ方程式を、\textbf{微分方程式}といいます。
微分方程式の解法第一号は、今回使える\textbf{変数分離法}です。
$g-\frac{k}{m}v\neq0$の範囲で、
\begin{align*}
  \frac{dv}{dt}=g-\frac{k}{m}v &\iff \frac{1}{g-kv/m}\frac{dv}{dt}=1
\end{align*}
この両辺を$t$で積分するのですが、右辺は簡単に$t+C$になりますね。
左辺は、合成関数の積分を使って
\begin{align*}
  \int \frac{1}{g-kv/m}\frac{dv}{dt}dt &= \int \frac{1}{g-kv/m}dv\\
  &=-\frac{m}{k}\log(g-kv/m)
\end{align*}
です。僕は暗算でこれくらいはできちゃいますが、どうやっているかというと「大体こういう感じになるだろうな」と目星をつけて、あとは適当に係数を決めてるだけです。

もとい。まとめて
\[
  -\frac{m}{k}\log(g-kv/m)=t+C \iff  g-\frac{k}{m}v(t)=A\exp\left(-\frac{kt}{m}\right) \iff v(t)=\frac{mg}{k}-\frac{mA}{k}\exp\left(-\frac{kt}{m}\right)
\]
と求めることができます。
積分定数$A$は、例えば自由落下するボールを「静かに落とした」という初期条件$v(0)=0$を課すとすると、
\[
  0=\frac{mg}{k}-\frac{mA}{k} \quad \therefore \; A=g
\]
と求まるので、空気抵抗付きの自由落下ボールの運動は
\[
  v(t)=\frac{mg}{k}\left\{1-\exp\left(-\frac{kt}{m}\right)\right\}
\]
で記述できるというわけです。

$t$が十分大きいとき、$e^{-\alpha t}$はかなり急速に0に収束します。
自由落下ボールの速度も、十分な時間を掛ければ$v=mg/k$に急速に近づいていきます。
例えば雨粒が雲の高さから空気抵抗を無視して地表に落ちてくると、とんでもないスピードで落ちてきます。
空気抵抗のおかげで、人間のたんぱく質を貫かない程度の勢いにとどめてくれているんですね。



\section{ばね運動と三角関数}

\subsection{ばね運動の運動方程式}

次のような運動方程式を考えます。
\begin{example}
  $\gamma:I\to\mathbb{R}$に対して、
  \begin{itemize}
    \item $m\ddot{\gamma}=-k\gamma$ ($k>0$)
    \item $\gamma(0)=1$
    \item $\dot\gamma(0)=0$ ($\omega>0$)
  \end{itemize}
  これを満たす$\gamma$を求めよ。
\end{example}

これは通常の運動方程式で、僕たちが解かねばならないタイプの方程式です。
物理的な状況としては、\textbf{ばねについた質点を、位置1で静かに手を離した}という状況です。
空気抵抗を無視しているので、ずっとビョンビョンするのです。
物理的には解$\gamma(t)=C(t)$が存在しそうですが、もう少しちゃんと議論する必要がありそうです。

まず簡単のため、運動方程式
\[
  \ddot{\gamma}(t)=-\gamma(t)
\]
を考えます。愚直にこれを積分すると、
\[
  \int_0^t\ddot{\gamma}(\tau)d\tau=-\int_0^t\gamma(\tau)d\tau
\]
となります。左辺は微分積分学の基本定理より
\[
  \int_0^t\ddot{\gamma}(\tau)d\tau\dot{\gamma}(t)-\dot\gamma(0)=\dot{\gamma}(t)
\]
ゆえ、
\[
  \dot{\gamma}(t)=-\int_0^t\gamma(\tau)d\tau
\]
さらに積分して、
\[
  \gamma(t)-1=-\int_0^t\int_0^{\tau}\gamma(\tau_1)d\tau_1d\tau
\]
従って、積分方程式と呼ばれるものの一つである次の方程式が得られます。
\[
  \gamma(t)=1-\int_0^t\int_0^{\tau}\gamma(\tau_1)d\tau_1d\tau
\]
これでは方程式が解けたのかどうかが分かりません。
筋が悪そうだからやめよう…などとは思わず、あきらめなかった男がピカールでした。
いったんこの右辺を
\[
  \gamma_1(t):=1-\int_0^t\int_0^{\tau}\gamma(\tau_1)d\tau_1d\tau
\]
とおきます。
さらに、$\gamma_1(t)$を積分方程式の右辺にぶちこんで、
\[
  \gamma_2(t):=1-\int_0^t\int_0^{\tau}\gamma_1(\tau_1)d\tau_1d\tau
\]
以下同様に、
\[
  \gamma_n(t):=1-\int_0^t\int_0^{\tau}\gamma_{n-1}(\tau_1)d\tau_1d\tau
\]
が得られます。

これについて考察を深めるために、具体的にいくつか計算してみましょう。
\begin{align*}
  \gamma_2(t)&=1-\int_0^t\int_0^{\tau}\gamma_1(\tau_1)d\tau_1d\tau\\
  &=1-\int_0^t\int_0^{\tau}\left[1-\int_0^{\tau_1}\int_0^{\tau_2}\gamma(\tau_3)d\tau_3d\tau_2\right]d\tau_1d\tau\\
  &=1-\frac{t^2}{2}+\int_0^t\int_0^{\tau}\int_0^{\tau_1}\int_0^{\tau_2}\gamma(\tau_3)d\tau_3d\tau_2d\tau_1d\tau
\end{align*}
\begin{align*}
  \gamma_3(t)&=1-\int_0^t\int_0^{\tau}\gamma_2(\tau_1)d\tau_1d\tau\\
  &=1-\int_0^t\int_0^\tau\left[1-\frac{\tau_1^2}{2}+\int_0^{\tau_1}\int_0^{\tau_2}\int_0^{\tau_3}\int_0^{\tau_4}\gamma(\tau_5)d\tau_5d\tau_4d\tau_3d\tau_2\right]d\tau_1d\tau\\
  &=1-\frac{t^2}{2}+\frac{t^4}{4!}-\int_0^t\int_0^\tau\int_0^{\tau_1}\int_0^{\tau_2}\int_0^{\tau_3}\int_0^{\tau_4}\gamma(\tau_5)d\tau_5d\tau_4d\tau_3d\tau_2d\tau_1d\tau
\end{align*}

目ざとい人には、これが$\cos$の級数表示になっていくことが理解できると思います。
しかし、今はまだ$\cos$が定義されていないことに注意してください。
僕たちはむしろ、$\ddot{\gamma}=-\gamma$の解として再定義しようとしているのです！
そのためには$\gamma_n$が、ある関数に収束することを証明しなければなりません。
このとき非常に扱いに困るのが、たくさん積分が並んでいる項
\[
  \int_0^t\int_0^\tau\int_0^{\tau_1}\int_0^{\tau_2}\int_0^{\tau_3}\int_0^{\tau_4}\gamma(\tau_5)d\tau_5d\tau_4d\tau_3d\tau_2d\tau_1d\tau
\]
などです。
\textbf{僕たちはこれを、思い切って無きものと思いましょう！}
どうやればそう思えるかというと、先ほど定義した関数列
\[
  \gamma_n(t):=1-\int_0^t\int_0^{\tau}\gamma_{n-1}(\tau_1)d\tau_1d\tau
\]
に対して、初期条件
\[
  \gamma_0(t)=0
\]
を課したものと考えます。
そして、極限
\[
  \gamma(t):=\lim_{n\to+\infty}\gamma_n(t)
\]
が存在し、これが微分可能かつ微分方程式
\[
  \ddot{\gamma}(t)=-\gamma(t)
\]
の解になっていることを確認できれば良いでしょう。

\subsection{$\gamma_n(t)$の"各点"収束}

$t\in\mathbb{R}$を固定すれば、関数列$\{\gamma_n(t)\}$はただの数列ですので、コーシーの収束判定法を使えます。
そこで、$\varepsilon>0$を任意の正の実数とします。
$|\gamma_n(t)-\gamma_m(t)|$を評価していきます。
$n\geq m>0$のとき、
\begin{align*}
  |\gamma_n(t)-\gamma_m(t)|&=\left|\int_0^t\int_0^\tau\gamma_{n-1}(\tau_1)d\tau_1d\tau-\int_0^t\int_0^\tau\gamma_{m-1}(\tau_1)d\tau_1d\tau\right|\\
  &=\left|\int_0^t\int_0^\tau\{\gamma_{n-1}(\tau_1)-\gamma_{m-1}(\tau_1)\}d\tau_1d\tau\right|\\
  &\leq\int_0^t\int_0^\tau|\gamma_{n-1}(\tau_1)-\gamma_{m-1}(\tau_1)|d\tau_1d\tau
\end{align*}
ここから再帰的に不等式を連ねることができますね。
いけるところまでやっちまいましょう。
\begin{align*}
  |\gamma_n(t)-\gamma_m(t)|&\leq\int_0^t\int_0^\tau|\gamma_{n-1}(\tau_1)-\gamma_{m-1}(\tau_1)|d\tau_1d\tau\\
  &\leq\int_0^t\int_0^\tau\int_0^{\tau_1}\int_0^{\tau_2}|\gamma_{n-2}(\tau_3)-\gamma_{m-2}(\tau_3)|d\tau_3d\tau_2d\tau_1d\tau\\
  &\cdots\\
  &\leq\int_0^t\int_0^\tau\int_0^{\tau_1}\cdots\int_0^{\tau_{2m-2}}|\gamma_{n-m}(\tau_{2m-1})-\gamma_0(\tau_{2m-1})|d\tau_{2m-1}\cdots d\tau_1d\tau
\end{align*}
$\gamma_0(t)=0$だったので、結局
\[
  |\gamma_n(t)-\gamma_m(t)|\leq\int_0^t\int_0^\tau\int_0^{\tau_1}\cdots\int_0^{\tau_{2m-2}}|\gamma_{n-m}(\tau_{2m-1})|d\tau_{2m-1}\cdots d\tau_1d\tau
\]
が分かりました。
